\section{LSS Foundation}

\label{sec:lss}

\subsection{Linear Secret Sharing}

All threshold protocols in \TChain{} rest on a common Linear Secret Sharing (\LSS{}) layer. A $(t, n)$-\LSS{} scheme distributes a secret $s$ into shares $(s_1, \ldots, s_n)$ via a degree-$(t{-}1)$ polynomial over $\mathbb{Z}_q$:

\begin{equation}
f(x) = s + a_1 x + a_2 x^2 + \cdots + a_{t-1} x^{t-1}, \quad s_i = f(i)
\end{equation}

Any $t$ shares reconstruct $s$ via Lagrange interpolation; fewer than $t$ shares reveal no information about $s$ (information-theoretic secrecy).

\subsection{The Composition Theorem}

The key insight enabling \TChain{}'s multi-protocol architecture is that protocols built on compatible \LSS{} instances can share the same access structure and secret shares.

\begin{theorem}[LSS Composition]
\label{thm:composition}
Let $\Pi_1, \ldots, \Pi_m$ be threshold cryptographic protocols, each requiring a $(t, n)$-\LSS{} over the same finite field $\mathbb{F}_q$. If each $\Pi_j$ uses shares of the form $s_i^{(j)} = f^{(j)}(i)$ for independent polynomials $f^{(j)}$ but the same evaluation points $\{1, \ldots, n\}$, then a single keygen ceremony can produce correlated shares for all protocols simultaneously via a vector polynomial:
\[
\mathbf{f}(x) = (f^{(1)}(x), \ldots, f^{(m)}(x))
\]
The resulting combined scheme preserves the $t$-privacy and $t$-correctness guarantees of each individual protocol.
\end{theorem}

This theorem is formally verified in \texttt{Crypto/Threshold/Composition.lean} and underpins the following protocol combinations in \TChain{}:

\begin{itemize}
\item \textbf{FROST}: Schnorr threshold signatures (Ed25519, Ristretto)
\item \textbf{CMP} (CGGMP21): ECDSA threshold signatures (secp256k1, secp256r1)
\item \textbf{TFHE}: Threshold fully-homomorphic encryption decryption
\item \textbf{Ringtail}: Ring signature scheme with threshold key generation
\item \textbf{BLS}: BLS threshold signatures for beacon chain compatibility
\end{itemize}

All five protocols derive their shares from a single DKG ceremony, reducing validator coordination overhead from $O(m)$ independent DKGs to a single vectorized DKG.

%% --------------------------------------------------------------------------
